\subsection{Square-root wake-ups from the two telescopes}
\label{sec:square-root-wakeups}

The diagonal-loss and correction-energy ledgers should not be charged
separately.  Their product is the actual uncertainty radius of a delayed
response.  This gives a square-root scheduling law before any particular
dynamic Schur representation is chosen.

\begin{lemma}[Two-ledger square-root wake-up bound]
\label{lem:two-ledger-square-root-wakeup}
Use the exact nested corrections and normalized frontier bases of
\cref{thm:causal-orthogonal-response-sketches}.  Fix a vertex $v$ throughout
an interval of events on which it remains exterior.  Partition those events
into disjoint consecutive epochs $I_1,\ldots,I_K$, and put
\begin{equation}
 A_r(v):=\sum_{j\in I_r}\lambda_{jv}^2,
 \qquad
 E_r:=\sum_{j\in I_r}\norm{\Delta^j}_Q^2.
 \label{eq:wakeup-two-ledgers}
\end{equation}
Then the normalized exterior-demand change accumulated at $v$ in epoch $r$
satisfies
\begin{equation}
 \frac{1}{\sqrt{d_v}}
 \left|e_v^TQ\sum_{j\in I_r}\Delta^j\right|
 \leq\sqrt{A_r(v)E_r}.
 \label{eq:wakeup-product-interval}
\end{equation}
Suppose epoch $r$ ends with a certified wake-up only after this product
radius reaches a declared positive margin $\mu_r$, so that
$\sqrt{A_r(v)E_r}\geq\mu_r$.  If
\begin{equation*}
 \mathcal E:=\sum_j\norm{\Delta^j}_Q^2,
 \qquad
 q_{\rm off}:=\frac{1-\alpha}{2},
\end{equation*}
then
\begin{equation}
 \sum_{r=1}^K\mu_r
 \leq
 \sqrt{\frac{q_{\rm off}}{d_v}\,\mathcal E}.
 \label{eq:wakeup-margin-budget}
\end{equation}
In particular, if every completed epoch uses margin at least $\mu>0$, then
\begin{equation}
 K\leq
 \frac{1}{\mu}
 \sqrt{\frac{q_{\rm off}}{d_v}\,\mathcal E}.
 \label{eq:wakeup-count}
\end{equation}

For a unit-mass seed and a trace starting from the zero face, the RPPR
objective gives $\mathcal E\leq\alpha$.  Hence at the finite-band scale
$\mu=c\alpha\tau$ for any fixed $c>0$,
\begin{equation}
 K=O\!\left(\frac{1}{\tau\sqrt{\alpha d_v}}\right).
 \label{eq:wakeup-product-scale}
\end{equation}
\end{lemma}

\begin{proof}
For each $j$, write $\Delta^j=U_ja_j$.  The normalized bases are mutually
$Q$-orthonormal, so
$\norm{a_j}_2=\norm{\Delta^j}_Q$.  Stack the row vectors
$e_v^TQU_j/\sqrt{d_v}$ and the coefficient vectors $a_j$ over
$j\in I_r$.  Ordinary Cauchy--Schwarz gives
\cref{eq:wakeup-product-interval}.

The epochs are disjoint.  A second Cauchy--Schwarz inequality and the wake-up
rule give
\begin{equation*}
 \sum_{r=1}^K\mu_r
 \leq\sum_{r=1}^K\sqrt{A_r(v)E_r}
 \leq
 \sqrt{\left(\sum_rA_r(v)\right)
             \left(\sum_rE_r\right)}.
\end{equation*}
By \cref{thm:source-leverage-diagonal-loss}, the first factor is the
normalized terminal-Schur diagonal loss over a subinterval of the exterior
lifetime of $v$ and is at most $q_{\rm off}/d_v$.  The second factor is at
most $\mathcal E$.  This proves
\cref{eq:wakeup-margin-budget,eq:wakeup-count}.

It remains to record the unit-seed energy budget.  The unregularized smooth
quadratic has minimum value
$-\tfrac12 b^TQ^{-1}b$.  Since
$Q^{-1}\preceq\alpha^{-1}I$ and
$b=\alpha D^{-1/2}s$ with $s\geq0$, $\one^Ts=1$, and $d_v\geq1$,
\begin{equation*}
 b^TQ^{-1}b
 \leq\alpha^{-1}\norm{b}_2^2
 \leq\alpha.
\end{equation*}
The nonnegative $\ell_1$ term can only raise the RPPR objective.  Therefore
the total objective drop from zero is at most $\alpha/2$, while
\cref{eq:orthogonal-correction-telescope} identifies twice that drop with
$\mathcal E$.  Thus $\mathcal E\leq\alpha$, and substituting
$\mu=c\alpha\tau$ proves \cref{eq:wakeup-product-scale}.
\end{proof}

The same argument applies to a fixed hierarchy node $C$ that remains
exterior: replace $A_r(v)$ by
$A_r(C):=\sum_{v\in C}\sum_{j\in I_r}\lambda_{jv}^2$ and the scalar response
by its $D_C^{-1/2}$-normalized vector.  The total first ledger is exactly the
aggregate Schur-diagonal loss in
\cref{eq:cumulative-sketch-diagonal-loss}.  Thus a hierarchy should not be
updated after every batch.  Each node should sleep while both its accumulated
diagonal loss and its accumulated orthogonal energy remain small, and wake
only when their product reaches the node's current pruning margin.

This lemma is a genuine square-root gain, but not yet the desired global work
theorem.  Summing the per-label bound naively can still pay for too many live
boundary labels, and maintaining the diagonal-loss ledger itself is the open
operation.  The useful conclusion is narrower: an implementation that
publishes every additive $\alpha\tau$ change discards an available
Cauchy--Schwarz gain.  The correct scheduler is multiplicative in response
leverage and correction energy.
